\chapter{Arquitectura de Datos}\label{cap:arquitectura}
\section{Introducción}
En este capítulo se detalla cómo TraDeck utiliza la red Ethereum y sistemas descentralizados para la persistencia y modelado de sus activos.

\section{El Ledger de Ethereum como Sistema de Persistencia}
La blockchain de Ethereum actúa como nuestra base de datos distribuida. A diferencia de un sistema SQL, la persistencia es inmutable y cada cambio de estado requiere un consenso en la red.

\section{Modelado de Activos: Estándares ERC-20 y ERC-721}
La arquitectura de TraDeck se apoya en dos estándares fundamentales de Ethereum para separar la lógica
financiera de la lógica de propiedad.

\subsection{Capa Financiera: ERC-20 (TraDeckCoin)}
Para gestionar la economía interna de la plataforma, se implementa un token ERC-20 \textbf{\cite{openzeppelin_erc20}} llamado TraDeckCoin (TDC) que
funciona como la moneda de intercambio dentro del mercado. Este token es fungible, lo que significa que cada unidad 
de TDC es idéntica a otra, permitiendo transacciones fluidas y eficientes.

\begin{itemize}
    \item \textbf{Estabilidad:} Permite fijar precios sin la volatilidad inherente a las criptomonedas tradicionales.
    \item \textbf{Utilidad:} Facilita la compra, venta e intercambio de cartas dentro de la plataforma sin depender 
    de monedas externas. 
\end{itemize}

\subsection{Capa de Propiedad: ERC-721 (TraDeckNFT)}
El estándar \textbf{ERC-721} \cite{openzeppelin_erc721} permite representar cada carta como un objeto 
único e indivisible. En nuestra arquitectura, el NFT actúa como la clave primaria y el certificado de 
propiedad irrefutable de cada carta coleccionable.

\subsubsection{Estructura de Datos y Mappings}
A diferencia de los sistemas de gestión de bases de datos tradicionales que usan tablas relacionales 
y claves foráneas, el ERC-721 asocia la propiedad mediante diccionarios de datos nativos de la 
blockchain (\verb|mapping|).

A través de funciones estándar como \verb|ownerOf(tokenId)|, el contrato inteligente actúa como nuestro motor 
de base de datos, permitiéndonos consultar qué dirección criptográfica (usuario) posee un ID de carta 
específico de forma instantánea y segura.

\subsubsection{Estrategia de Almacenamiento y Costes}
Almacenar grandes cantidades de datos en la blockchain, como imágenes de alta calidad, tiene un coste de 
transacción alto y prohibitivo. La propia documentación de OpenZeppelin advierte sobre esto y sugiere el uso 
de sistemas como IPFS.

Por ello, aplicamos una división de datos:
\begin{itemize}
    \item \textbf{Datos On-Chain}. \\ Solo se persisten en la blockchain los datos considerados críticos 
    que garantizan la titularidad:
    \begin{itemize}
        \item \textbf{tokenId:} ID de la carta. 
        \item \textbf{owner:} Dirección del propietario.
        \item \textbf{tokenURI:} Un puntero o enlace al alojamiento externo.
    \end{itemize}

    \item \textbf{Datos Off-Chain}. \\ El \texttt{tokenURI} se diseña para que resuelva a un documento JSON alojado en la 
    red descentralizada IPFS. Este archivo contiene los datos pesados de la carta: 

    \begin{lstlisting}[language=json, caption=Ejemplo de tokenURI apuntando a IPFS]
{
    "name": "Charizard",
    "description": "Carta holografica primera edicion",
    "image": "ipfs://<CID-de-la-imagen>"
}
    \end{lstlisting}

\end{itemize}

\subsection{¿Por qué no utilizar el ERC-1155?}
Durante la fase de investigación, se evaluó el uso del estándar ERC-1155, que permite la gestión de mútliples tipos de 
tokens (fungibles y no fungibles) dentro de un mismo contrato. Sin embargo, se descartó por los siguientes motivos:

\begin{itemize}
    \item \textbf{Unicidad Estricta:} TraDeck certifica cartas coleccionables donde cada ejemplar es único, con un 
    estado físico y número de serie específicos. El ERC-721 garantiza una trazabilidad 1 a 1 estricta, mientras que 
    el ERC-1155 está optimizado para 'ediciones' o lotes de objetos idénticos, lo que no se ajustaba a nuestro modelo de negocio. 

    \item \textbf{Responsabilidad Única:} Separar la lógica en dos contratos, ERC-20 y ERC-721, sigue los principios 
    de diseño modular, facilitando la seguridad y la auditabilidad. El patrón Escrow, que es crucial para garantizar
    transacciones seguras, se implementa de manera más clara y efectiva con esta separación, evitando la complejidad 
    adicional que introduciría el ERC-1155.
\end{itemize}

\section{Integridad Referencial: Almacenamiento Off-chain mediante IPFS}

La decisión de almacenar los metadatos y los activos multimedia fuera de la blockchain (off-chain) resuelve el problema de los costes de persistencia pero introduce otro nuevo desafío arquitectónico: \textbf{¿Cómo podemos garantizar que los datos guardados off-chain, externos, no son alterados, perdiendo así la integridad de la base de datos distribuida?}

Si utilizáramos un servidor centralizado tradicional (como Amazon S3 o incluso Google Drive) mediante un enlace HTTP estándar, estaríamos utilizando un \textit{direccionamiento por ubicación} (Location-Addressing), lo cual permite a un administrador o cualquier intruso no autorizado alterar el archivo JSON o la imagen en el servidor sin que cambie la URL: el Smart Contract apuntaría a la misma dirección pero el activo digital habría sido falsificado.

Para solucionar este problema de integridad, TraDeck utiliza \textbf{IPFS (InterPlanetary File System)}, una red de almacenamiento descentralizada basada en \textit{direccionamiento por contenido} (Content-Addressing):

\begin{itemize}
	\item \textbf{Identificador de Contenido (CID):} En lugar de usar una ruta estática, IPFS procesa el archivo (la imagen o el JSON) a través de una función hash criptográfica, generando un identificador único (CID).
	\item \textbf{Garantía de Inmutabilidad:} El \texttt{tokenURI} que se persiste en la blockchain es, en realidad, este CID. Si cualquier nodo o usuario intentara modificar un solo bit de la imagen o de la descripción de la carta, el hash resultante cambiaría, sería completamente distinto. 
	\item \textbf{Verificación P2P:} Como la blockchain guarda el CID original de forma inmutable, cualquier alteración en el almacenamiento off-chain rompería el enlace en la cadena, el enlace criptográfico. Esto hace que la base de datos híbrida (Blockchain + IPFS) sea a prueba de manipulaciones sin requerir confianza en un servidor central.
\end{itemize}

Mediante esta arquitectura, logramos un equilibrio: la eficiencia económica de no saturar el ledger de Ethereum, manteniendo la robustez transaccional e inmutabilidad estricta exigida en un sistema de bases de datos de alta seguridad.

